% ============================================================
%  Resumen — Extensiones de Machine Learning: Aprendizaje por Refuerzo
% ============================================================

\begin{abstract}
\noindent
Este trabajo estudia el Aprendizaje por Refuerzo (RL) de forma progresiva, partiendo del paradigma más simple de toma de decisiones secuencial hasta los métodos de aproximación profunda con redes neuronales, a través de tres bloques de contenido interconectados.

\medskip
\noindent
\textbf{Capítulo 1 — Problema del Bandido de $k$-brazos.}
Se introduce el dilema fundamental del RL: el balance entre \textit{exploración} y \textit{explotación}. Un agente interactúa con un entorno estocástico de $k$ palancas, cada una con una distribución de recompensa desconocida, buscando maximizar la ganancia acumulada en un horizonte finito de $T = 300$ pasos. Se implementan y comparan empíricamente tres familias algorítmicas bajo $N = 500$ ejecuciones y distribuciones de recompensa Normal, Binomial y Bernoulli: (i)~\textbf{$\epsilon$-Greedy} y su variante con decaimiento $\epsilon$-Decay, que combinan explotación determinista con exploración aleatoria; (ii)~\textbf{UCB} (UCB1, UCB2, UCB1-Tuned), que racionalizan la exploración mediante límites de confianza superior sobre la incertidumbre de cada brazo; y (iii)~\textbf{Ascenso de Gradiente} (Softmax y Preferencias), que aprenden una distribución de selección sobre las acciones. Los resultados muestran que no existe un algoritmo universalmente óptimo: UCB2 sobresale en entornos Normales, mientras que el Gradiente de Preferencias ($\alpha = 0.5$) domina en el régimen binario Bernoulli con un arrepentimiento acumulado prácticamente nulo ($\text{Regret} \approx 0$).

\medskip
\noindent
\textbf{Capítulo 2 — Entornos Complejos: Métodos Tabulares.}
Se formula el problema de control óptimo como un Proceso de Decisión de Markov (MDP) $\langle \mathcal{S}, \mathcal{A}, \mathcal{P}, \mathcal{R}, \gamma \rangle$ y se evalúa la familia de métodos tabulares model-free sobre el entorno \texttt{CliffWalking-v0}. Se implementan Monte Carlo On/Off-Policy, SARSA, Expected SARSA, Q-Learning, Double Q-Learning y variantes de $n$-pasos. Se diseña una arquitectura modular que desacopla el agente, el algoritmo (\textit{Learner}) y el entorno (Gymnasium), garantizando reproducibilidad con $N = 10$ ejecuciones de 400 episodios ($\gamma = 0.95$, $\alpha = 0.5$). Los experimentos evidencian la superioridad de Expected SARSA ($\approx{-}15$ puntos de penalización) frente a Q-Learning ($\approx{-}56$) gracias a su robustez ante tasas de aprendizaje elevadas, y demuestran la ineficacia práctica de Monte Carlo en entornos con penalizaciones severas. Double Q-Learning resuelve el \textit{Maximization Bias} a costa de una convergencia más lenta. El análisis concluye con el límite intrínseco de los métodos tabulares: la \textit{maldición de la dimensionalidad} hace inviable la representación matricial $\mathbf{Q} \in \mathbb{R}^{|\mathcal{S}| \times |\mathcal{A}|}$ en espacios de estados continuos.

\medskip
\noindent
\textbf{Capítulo 3 — Entornos Complejos: Métodos Aproximados.}
Para superar la barrera de escalabilidad, se substituye la tabla $\mathbf{Q}$ por una función paramétrica $\hat{q}(\cdot\,;\,\boldsymbol{\theta})$ y se evalúan tres aproximadores de complejidad creciente sobre \texttt{CartPole-v1} ($\mathcal{S} \subseteq \mathbb{R}^4$) con $N = 10$ ejecuciones de 400 episodios: (i)~\textbf{SARSA Semi-Gradiente} con Tile Coding ($d = 16\,384$ características binarias, 32\,768 parámetros), cuya aproximación lineal limita la expresividad; (ii)~\textbf{DQN} (red $4 \to 64 \to 2$ con ReLU, Adam $\eta = 10^{-3}$, Experience Replay y red target), que combina la no linealidad neuronal con estabilización del proceso de Bellman; y (iii)~\textbf{Double DQN}, extensión de DQN que desacopla selección y evaluación de la acción óptima para corregir el sesgo de sobreestimación. Los resultados ($\overline{R}_{[-50:]}$: SARSA~69, DQN~218, Double DQN~189) confirman que la capacidad expresiva del aproximador determina la calidad de la política, que el Experience Replay es condición crítica de estabilidad, y que la ventaja de Double DQN sobre DQN solo emerge en horizontes de entrenamiento superiores a $T_\text{ep} > 1\,000$ episodios.

\medskip
\noindent
En conjunto, los tres capítulos trazan una progresión coherente desde la exploración-explotación en un único estado hasta el control de sistemas dinámicos continuos mediante redes neuronales profundas, ilustrando los compromisos fundamentales del RL entre \textbf{capacidad expresiva}, \textbf{eficiencia muestral} y \textbf{estabilidad del aprendizaje}.
\end{abstract}
